% ======================================================================
% Chapter 5: Elementary Processes of Quantum Electrodynamics
% (book pages 131-173; PDF pages 154-197)
% ======================================================================
\chapter{量子电动力学的基本过程}

在经历三章冗长的形式主义之后，我们终于准备好进行一些真正的相对论计算，开始推导量子电动力学的预言。首先我们将回到第一章中考虑的过程：电子-正电子对湮灭为一对更重的费米子。我们将在接下来的三节中极其详细地研究这一典范过程，然后在第5.4节和第5.5节中进行一些更简单的 QED 计算。章末的习题涉及若干额外的 QED 过程。更完整的 QED 综述可见于 Jauch 和 Rohrlich (1976) 以及 Berestetskii、Lifshitz 和 Pitaevskii (1982) 的著作。

\section{$e^{+}e^{-}\to\mu^{+}\mu^{-}$：引言}

反应 $e^{+}e^{-}\to\mu^{+}\mu^{-}$ 是所有 QED 过程中最简单的一个，但也是高能物理中最重要的过程之一。它是理解 $e^{+}e^{-}$ 对撞机中所有反应的基础，实际上还被用来校准这类机器。相关的过程 $e^{+}e^{-}\to q\bar{q}$（夸克-反夸克对）在确定基本粒子性质方面极为有用。

在本节中，我们将计算最低阶的 $e^{+}e^{-}\to\mu^{+}\mu^{-}$ 未极化截面。在第一章中，我们在所有费米子无质量的极限下，用初等论证猜出了答案（式~\eqref{eq:1.8}）。现在我们放宽这一限制，在计算中保留缪子质量。同时保留电子质量也很容易，但没有必要，因为比值 $m_e/m_\mu \approx 1/200$ 远小于忽略微扰级数中高阶项所引入的相对误差。

利用第4.8节的费曼规则，我们可以立即画出该过程的费曼图并写出其振幅：
\begin{equation}
  i\mathcal{M} = \bar{v}^{s'}(p')(-ie\gamma^\mu)u^s(p)\
    \frac{-ig_{\mu\nu}}{q^2}\ \bar{u}^{r}(k)(-ie\gamma^\nu)v^{r'}(k').
\end{equation}
稍加整理并隐去自旋上标，我们得到
\begin{equation}
  i\mathcal{M}\bigl( e^{-}(p)e^{+}(p') \to \mu^{-}(k)\mu^{+}(k') \bigr)
    = \frac{(-ie)^2}{q^2}\ \bigl( \bar{v}(p')\gamma^\mu u(p) \bigr)\
      \bigl( \bar{u}(k)\gamma_\mu v(k') \bigr) .
  \label{eq:5.1}
\end{equation}
振幅 $\mathcal{M}$ 的这个结果虽然简单，但还不够直观。为了计算微分截面，我们需要 $|\mathcal{M}|^2$ 的表达式，因此必须求出 $\mathcal{M}$ 的复共轭。像 $\bar{v}\gamma^\mu u$ 这样的双旋量乘积可以按如下方式取复共轭：
\begin{equation}
  \bigl( \bar{v}\gamma^\mu u \bigr)^{*}
    = u^{\dagger}\gamma^{\mu\dagger}(\gamma^0)^{\dagger}v
    = \bar{u}\gamma^\mu v .
\end{equation}
（这是``bar''记号的另一个优点。）因此，矩阵元的平方为
\begin{equation}
  |\mathcal{M}|^2 = \frac{e^4}{q^4}\ \bigl( \bar{v}(p')\gamma^\mu u(p)
    \bar{u}(p)\gamma^\nu v(p') \bigr)\
    \bigl( \bar{u}(k)\gamma_\mu v(k')\ \bar{v}(k')\gamma_\nu u(k) \bigr).
  \label{eq:5.2}
\end{equation}
在这一点上，我们仍然可以自由地指定任意特定的旋量 $u^s(p)$、$v^{s'}(p')$ 等，对应于费米子的任意所需自旋态。然而，在实际实验中，保持对自旋态的控制是困难的（尽管并非不可能）；人们需要用极化的材料来制备初态，和/或用依赖于自旋的多重散射来分析末态。在大多数实验中，电子束和正电子束都是未极化的，因此测得的截面是电子和正电子自旋 $s$ 与 $s'$ 的平均。缪子探测器通常对极化不敏感，因此测得的截面是对缪子自旋 $r$ 与 $r'$ 的求和。

当我们丢弃自旋信息时，$|\mathcal{M}|^2$ 的表达式会大大简化。我们希望计算
\begin{equation}
  \frac{1}{4}\sum_{s,s'}\sum_{r,r'}|\mathcal{M}|^2 .
\end{equation}
自旋求和可以用第3.3节的完备性关系来完成：
\begin{equation}
  \sum_s u^s(p)\bar{u}^s(p) = \slashed{p} + m;\qquad
  \sum_s v^s(p)\bar{v}^s(p) = \slashed{p} - m .
  \label{eq:5.3}
\end{equation}
处理 \eqref{eq:5.2} 的前半部分，并写出旋量指标以便自由地把 $v$ 移到 $\bar{v}$ 旁边，我们有
\begin{equation}
  \begin{aligned}
    \sum_{s,s'} (\slashed{p}' - m)_{da}\gamma^\mu_{ab}(\slashed{p} + m)_{bc}
      \gamma^\nu_{cd} &= \mathrm{tr}\bigl[
        (\slashed{p}' - m)\gamma^\mu(\slashed{p} + m)\gamma^\nu \bigr].
  \end{aligned}
\end{equation}
用同样的方法计算 \eqref{eq:5.2} 的后半部分，我们得到所需的简化：
\begin{equation}
  \frac{1}{4}\sum_{\text{spins}}|\mathcal{M}|^2
    = \frac{e^4}{4q^4}\ \mathrm{tr}\bigl[
      (\slashed{p}' - m_e)\gamma^\mu(\slashed{p} + m_e)\gamma^\nu \bigr]\
      \mathrm{tr}\bigl[ (\slashed{k} + m_\mu)\gamma_\mu(\slashed{k}' - m_\mu)
        \gamma_\nu \bigr] .
  \label{eq:5.4}
\end{equation}
旋量 $u$ 和 $v$ 已经消失，留下一个用 $\gamma$ 矩阵表示的简洁得多的表达式。这个技巧非常普遍：任何涉及外线费米子的 QED 振幅，在取平方并对自旋求和或平均后，都可以用这种方式转化为 Dirac 矩阵乘积的迹。

\subsection{求迹技巧}

如果这些迹必须费力地用蛮力计算，那么最后这一步就谈不上什么改进。但 Feynman 发现，利用 $\gamma$ 矩阵的代数性质可以很容易地算出它们。由于这类迹的求值在 QED 计算中如此频繁地出现，值得我们停下来一劳永逸地系统处理这个问题。

我们希望求 $n$ 个伽马矩阵乘积的迹，其中 $n = 0, 1, 2, \ldots$。（对于当前问题，我们需要 $n = 2, 3, 4$。）$n=0$ 的情形相当容易：$\mathrm{tr}\,1 = 4$。一个 $\gamma$ 矩阵的迹也很容易。根据手征表象中矩阵的显式形式，我们有
\begin{equation}
  \mathrm{tr}\,\gamma^\mu = 0 .
\end{equation}
用一种更抽象的方法来证明这个结果是有用的，它可以推广到任意奇数个 $\gamma$ 矩阵的情形：
\begin{equation}
  \begin{aligned}
    \mathrm{tr}\,\gamma^\mu &= \mathrm{tr}\,\gamma^5\gamma^5\gamma^\mu
      &&\text{因为 } (\gamma^5)^2 = 1\\
    &= -\mathrm{tr}\,\gamma^5\gamma^\mu\gamma^5 &&\text{因为 }
      \{\gamma^\mu, \gamma^5\} = 0\\
    &= -\mathrm{tr}\,\gamma^5\gamma^5\gamma^\mu
      &&\text{利用迹的循环性质}\\
    &= -\mathrm{tr}\,\gamma^\mu .
  \end{aligned}
\end{equation}
由于 $\gamma^\mu$ 的迹等于其自身取负，它必然为零。对于 $n$ 个 $\gamma$ 矩阵，我们在第二步中会得到 $n$ 个负号（当我们把第二个 $\gamma^5$ 一直移到最右边时），因此当 $n$ 为奇数时迹必定为零。

为了求两个 $\gamma$ 矩阵的迹，我们再次使用反对易性质和迹的循环性质：
\begin{equation}
  \begin{aligned}
    \mathrm{tr}\,\gamma^\mu\gamma^\nu
      &= \mathrm{tr}(2g^{\mu\nu}\cdot 1 - \gamma^\nu\gamma^\mu)
      &&\text{（反对易）}\\
    &= 8g^{\mu\nu} - \mathrm{tr}\,\gamma^\mu\gamma^\nu
      &&\text{（循环性）}
  \end{aligned}
\end{equation}
因此 $\mathrm{tr}\,\gamma^\mu\gamma^\nu = 4g^{\mu\nu}$。任何偶数个 $\gamma$ 矩阵的迹都可以用同样的方法求值：把第一个 $\gamma$ 矩阵一直反对易到最右边，然后循环回最左边。因此，对于四个 $\gamma$ 矩阵的迹，我们有
\begin{equation}
  \begin{aligned}
    \mathrm{tr}\,(\gamma^\mu\gamma^\nu\gamma^\rho\gamma^\sigma)
    &= \mathrm{tr}(2g^{\mu\nu}\gamma^\rho\gamma^\sigma
      - \gamma^\nu\gamma^\mu\gamma^\rho\gamma^\sigma)\\
    &= \mathrm{tr}(2g^{\mu\nu}\gamma^\rho\gamma^\sigma
      - 2g^{\mu\rho}\gamma^\nu\gamma^\sigma
      + 2g^{\mu\sigma}\gamma^\nu\gamma^\rho
      - \gamma^\nu\gamma^\rho\gamma^\sigma\gamma^\mu) .
  \end{aligned}
\end{equation}
对最后一项使用循环性质并将其移到左边，我们得到
\begin{equation}
  \mathrm{tr}(\gamma^\mu\gamma^\nu\gamma^\rho\gamma^\sigma)
    = g^{\mu\nu}\mathrm{tr}\,\gamma^\rho\gamma^\sigma
      - g^{\mu\rho}\mathrm{tr}\,\gamma^\nu\gamma^\sigma
      + g^{\mu\sigma}\mathrm{tr}\,\gamma^\nu\gamma^\rho
    = 4\bigl( g^{\mu\nu}g^{\rho\sigma} - g^{\mu\rho}g^{\nu\sigma}
      + g^{\mu\sigma}g^{\nu\rho} \bigr) .
\end{equation}
用这种方式，人们总能将一个 $n$ 个 $\gamma$ 矩阵的迹约化为 $(n-2)$ 个 $\gamma$ 矩阵的迹之和。$n=6$ 的情形不难算出，但它有十五项（把六个指标两两配对以构成 $g^{\mu\nu}g^{\rho\sigma}g^{\alpha\beta}$ 形式的项的方式数）。幸运的是，本书中我们用不到它。（如果你确实需要求这样的复杂迹，学习使用几个能够对 Dirac 矩阵进行符号运算的计算机程序之一可能会更容易。）

从第5.2节开始，我们将经常需要求涉及 $\gamma^5$ 的迹。由于 $\gamma^5 = i\gamma^0\gamma^1\gamma^2\gamma^3$，$\gamma^5$ 与任意奇数个其他 $\gamma$ 矩阵乘积的迹为零。同样容易证明 $\gamma^5$ 自身的迹为零：
\begin{equation}
  \mathrm{tr}\,\gamma^5 = \mathrm{tr}(\gamma^0\gamma^0\gamma^5)
    = -\mathrm{tr}(\gamma^0\gamma^5\gamma^0)
    = -\mathrm{tr}(\gamma^0\gamma^0\gamma^5) = -\mathrm{tr}\,\gamma^5 .
\end{equation}
同样的技巧也适用于 $\mathrm{tr}(\gamma^\mu\gamma^\nu\gamma^5)$，只要我们插入两个因子 $\gamma^\alpha$，其中 $\alpha$ 不同于 $\mu$ 和 $\nu$。涉及 $\gamma^5$ 的第一个非零迹包含四个其他 $\gamma$ 矩阵。在这种情况下，除非每个 $\gamma$ 矩阵都出现，否则该技巧仍然有效，因此
\begin{equation}
  \mathrm{tr}(\gamma^\mu\gamma^\nu\gamma^\rho\gamma^\sigma\gamma^5) = 0
  \qquad\text{除非 } (\mu\nu\rho\sigma) \text{ 是 }
    (0123) \text{ 的某个置换}.
\end{equation}
从反对易规则还可推出，交换任意两个指标只会改变迹的符号，因此 $\mathrm{tr}(\gamma^\mu\gamma^\nu\gamma^\rho\gamma^\sigma\gamma^5)$ 必定正比于 $\epsilon^{\mu\nu\rho\sigma}$。总的常数经计算为 $-4i$，你只需代入 $(\mu\nu\rho\sigma) = (0123)$ 即可容易地验证。

下面是求迹定理的总结，方便查阅：
\begin{equation}
  \begin{aligned}
    \mathrm{tr}(1) &= 4\\
    \mathrm{tr}(\text{任意奇数个 }\gamma\text{ 的乘积}) &= 0\\
    \mathrm{tr}(\gamma^\mu\gamma^\nu) &= 4g^{\mu\nu}\\
    \mathrm{tr}(\gamma^\mu\gamma^\nu\gamma^\rho\gamma^\sigma)
      &= 4\bigl( g^{\mu\nu}g^{\rho\sigma} - g^{\mu\rho}g^{\nu\sigma}
        + g^{\mu\sigma}g^{\nu\rho} \bigr)\\
    \mathrm{tr}(\gamma^5) &= 0\\
    \mathrm{tr}(\gamma^\mu\gamma^\nu\gamma^5) &= 0\\
    \mathrm{tr}(\gamma^\mu\gamma^\nu\gamma^\rho\gamma^\sigma\gamma^5)
      &= -4i\epsilon^{\mu\nu\rho\sigma}
  \end{aligned}
  \label{eq:5.5}
\end{equation}
由最后一条公式得出的表达式可以借助下列恒等式加以简化：
\begin{equation}
  \epsilon_{\alpha\beta\gamma\delta}\epsilon^{\alpha\beta\gamma\delta}
    = -24;\qquad
  \epsilon_{\alpha\beta\mu\nu}\epsilon^{\alpha\beta\rho\sigma}
    = -2\bigl( \delta^\rho_\mu\delta^\sigma_\nu
      - \delta^\rho_\nu\delta^\sigma_\mu \bigr) .
  \label{eq:5.6}
\end{equation}
所有这些都可以先借助对称性论证，然后计算一个特例以确定总的常数来导出。

另一个有用的恒等式允许人们颠倒迹内所有 $\gamma$ 矩阵的顺序：
\begin{equation}
  \mathrm{tr}(\gamma^\mu\gamma^\nu\gamma^\rho\gamma^\sigma)
    = \mathrm{tr}(\gamma^\sigma\gamma^\rho\gamma^\nu\gamma^\mu) .
  \label{eq:5.7}
\end{equation}

我们现在可以完成 \eqref{eq:5.4} 的求值。舍弃电子质量（令 $m_e = 0$），我们得到
\begin{equation}
  \frac{1}{4}\sum_{\text{spins}}|\mathcal{M}|^2
    = \frac{e^4}{4q^4}\ \bigl[ 4\bigl( p'^{\mu}p^{\nu} - p'\cdot p\,g^{\mu\nu}
      + p'^{\nu}p^{\mu} \bigr) \bigr]\ \bigl[ 4\bigl( k_\mu k'_\nu
        - k\cdot k'\,g_{\mu\nu} + k_\nu k'_\mu \bigr) \bigr].
  \label{eq:5.8}
\end{equation}
缩并指标并利用 $q^2 = (p+p')^2 = 2p\cdot p'$，我们得到
\begin{equation}
  \frac{1}{4}\sum_{\text{spins}}|\mathcal{M}|^2
    = \frac{8e^4}{q^4}\ \bigl[ (p\cdot k)(p'\cdot k') + (p\cdot k')(p'\cdot k)
      + m_\mu^2\,p\cdot p' \bigr] .
  \label{eq:5.9}
\end{equation}
将结果特化到质心系，其中 $E$ 是电子（和正电子）的能量，$E_\mu$ 是缪子（和反缪子）的能量，我们有 $p = (E, \mathbf{p})$，$p' = (E, -\mathbf{p})$，$k = (E_\mu, \mathbf{k})$，$k' = (E_\mu, -\mathbf{k})$，且 $q^2 = 4E^2$。直接计算给出
\begin{equation}
  \frac{1}{4}\sum_{\text{spins}}|\mathcal{M}|^2
    = \frac{e^4}{4E^4}\ \bigl[ (E + E_\mu)^2 + (E\cos\theta + E_\mu)^2
      + m_\mu^2 \bigr] .
  \label{eq:5.10}
\end{equation}
截面随后由一般公式 \eqref{eq:4.82} 得出：
\begin{equation}
  \Bigl( \frac{d\sigma}{d\Omega} \Bigr)_{\mathrm{CM}}
    = \frac{1}{64\pi^2 E_{\mathrm{cm}}^2}\ \frac{1}{4}
      \sum_{\text{spins}}|\mathcal{M}|^2 .
  \label{eq:5.11}
\end{equation}
代入 \eqref{eq:5.10}，并利用 $E_{\mathrm{cm}} = 2E$，$E_\mu = E_{\mathrm{cm}}/2$ 进行化简，我们得到
\begin{equation}
  \Bigl( \frac{d\sigma}{d\Omega} \Bigr)_{\mathrm{CM}}
    = \frac{\alpha^2}{4E_{\mathrm{cm}}^2}\ \sqrt{1 - \frac{m_\mu^2}{E^2}}
      \Bigl[ 1 + \cos^2\theta + \frac{m_\mu^2}{E^2}\sin^2\theta \Bigr] ,
  \label{eq:5.12}
\end{equation}
其中 $\alpha = e^2/4\pi$。对 $\theta$ 和 $\phi$ 积分给出总截面，
\begin{equation}
  \sigma_{\mathrm{total}}
    = \frac{4\pi\alpha^2}{3E_{\mathrm{cm}}^2}\ \sqrt{1 - \frac{m_\mu^2}{E^2}}
      \Bigl[ 1 + \frac{m_\mu^2}{2E^2} \Bigr] .
  \label{eq:5.13}
\end{equation}
在高能极限 $E \gg m_\mu$ 下，这些公式约化为第一章中给出的结果：
\begin{equation}
  \frac{d\sigma}{d\Omega} \to \frac{\alpha^2}{4E_{\mathrm{cm}}^2}
    (1 + \cos^2\theta);\qquad
  \sigma_{\mathrm{total}} \to \frac{4\pi\alpha^2}{3E_{\mathrm{cm}}^2} .
  \label{eq:5.14}
\end{equation}
注意这些表达式具有截面的正确量纲。在高能极限下，$E_{\mathrm{cm}}$ 是问题中唯一具有量纲的量，因此量纲分析决定了 $\sigma_{\mathrm{total}} \sim E_{\mathrm{cm}}^{-2}$。由于我们从一开始就知道 $\sigma_{\mathrm{total}} \propto \alpha^2$，我们只需计算以得到因子 $4\pi/3$。

总截面公式 \eqref{eq:5.13} 在阈值附近的能量依赖关系如图~\ref{fig:5.1} 所示。当然，对于 $E_{\mathrm{cm}} < 2m_\mu$，截面为零。将实际曲线的形状与假设 $|\mathcal{M}|^2$ 不依赖于能量时得到的形状（即所有能量依赖均来自相空间因子 $|\mathbf{k}|/E$）进行比较是很有趣的。为了检验量子电动力学，实验必须能够分辨出对朴素相空间预言的偏离。$\mu$ 和 $\tau$ 轻子对产生的实验结果都证实这些粒子遵从 QED 的预言。图~\ref{fig:5.2} 将公式 \eqref{eq:5.13} 与 $\tau^{+}\tau^{-}$ 阈值的实验测量进行了比较。

\begin{figure}[t]
  \centering
  \includegraphics[width=0.6\textwidth]{fig5.1.png}
  \caption{$e^{+}e^{-}\to\mu^{+}\mu^{-}$ 总截面的能量依赖关系与``相空间''能量依赖关系的比较。}
  \label{fig:5.1}
\end{figure}

\begin{figure}[t]
  \centering
  \includegraphics[width=0.7\textwidth]{fig5.2.png}
  \caption{DELCO 合作组（W.~Bacino 等人，\emph{Phys.\ Rev.\ Lett.} \textbf{41}，13 (1978)）测量的 $\tau^{+}\tau^{-}$ 对产生阈值附近的截面比 $\sigma(e^{+}e^{-}\to\tau^{+}\tau^{-})/\sigma(e^{+}e^{-}\to\mu^{+}\mu^{-})$。只包含了 $\tau$ 衰变的一小部分，因此整体标度较小。曲线是对理论公式 \eqref{eq:5.13} 的拟合，并加上了小的能量无关本底。拟合给出 $m_\tau = 1782\frac{+2}{-2}$ MeV。}
  \label{fig:5.2}
\end{figure}

在进一步讨论我们的结果之前，让我们停下来总结一下我们是如何得到它的。该方法可以直接推广到其他 QED 过程未极化截面的计算。一般步骤如下：
\begin{enumerate}
  \item 画出所要求过程的费曼图。
  \item 用费曼规则写出振幅 $\mathcal{M}$。
  \item 将振幅平方并利用完备性关系 \eqref{eq:5.3} 对自旋求平均或求和。（对于末态含光子的过程，存在类似的完备性关系，将在第5.5节中导出。）
  \item 利用求迹定理 \eqref{eq:5.5} 求迹；合并各项并尽可能简化结果。
  \item 特化到某个特定的参考系，并在该参考系中画出运动学变量的示意图。将所有四动量矢量用适当选取的一组变量（如 $E$ 和 $\theta$）来表示。
  \item 将所得的 $|\mathcal{M}|^2$ 表达式代入截面公式 \eqref{eq:4.79}，并对未测量的相空间变量积分，以获得所需形式的微分截面。（在我们的情形中，这些积分是对受约束的动量 $\mathbf{k}'$ 和 $|\mathbf{k}|$ 进行的，已在式~\eqref{eq:4.84} 的推导中完成。）
\end{enumerate}
虽然其他计算（尤其是涉及圈图的那些）常常需要额外的技巧，但几乎每一个 QED 计算都会涉及这里概述的基本步骤。

\subsection{夸克-反夸克对的产生}

$e^{+}e^{-}\to\mu^{+}\mu^{-}$ 截面公式的渐近能量依赖关系为所有 $e^{+}e^{-}$ 湮灭截面设定了标度。一个特别重要的例子是下列过程的截面
\begin{equation}
  e^{+}e^{-} \to \text{强子},
\end{equation}
即产生任意数目强相互作用粒子的总截面。

按照我们目前对强相互作用的理解（由称为量子色动力学（QCD）的理论给出），所有强子都由称为夸克的 Dirac 费米子组成。夸克以各种类型出现，称为味，每一种都有自身的质量和电荷。夸克还携带一个附加的量子数——色，它取三个值之一。色充当 QCD 的``电荷''，我们将在第17章中讨论。

根据 QCD，以强子为末态的最简单的 $e^{+}e^{-}$ 过程是
\begin{equation}
  e^{+}e^{-} \to q\bar{q},
\end{equation}
即一个电子和一个正电子通过一个虚光子湮灭为夸克-反夸克对。在产生之后，夸克之间通过强力相互作用，产生更多的夸克对。最终夸克和反夸克结合形成若干介子和重子。

为了将我们关于缪子产生的结果推广到夸克的情形，我们必须做三处修改：
\begin{enumerate}
  \item 将缪子电荷 $e$ 替换为夸克电荷 $Q|e|$。
  \item 对每个夸克计数三次，每种颜色计一次。
  \item 计入所产生的夸克和反夸克的强相互作用效应。
\end{enumerate}
前两处修改很容易实现。对于第一处，只需要知道每种味夸克的质量和电荷。对于 $u$、$c$ 和 $t$ 夸克，我们有 $Q = 2/3$，而对于 $d$、$s$ 和 $b$ 夸克，我们有 $Q = -1/3$。截面公式正比于末态粒子电荷的平方，因此我们可以简单地在这些公式中的任何一个插入因子 $Q^2$，以获得产生任何特定种类夸克的截面。对颜色计数是必要的，因为实验只测量所有三种颜色产生的总截面。（实际探测到的强子是无色的。）无论如何，这种计数很简单：只需把结果乘以 3。

然而，如果你对强相互作用稍有所知，你可能会认为这一切都是个大玩笑。第三处修改肯定极其困难，并且会彻底改变 QED 的预言。令人惊讶的事实是，在高能极限下，强相互作用对夸克产生过程的影响可以完全忽略。正如我们将在第三部分中讨论的，强相互作用的唯一效应（在此极限下）就是把末态夸克装扮成成束的强子。这种简化归因于一个称为渐近自由的现象；它在确立量子色动力学为强相互作用的正确理论中起到了关键作用。

因此在极限下，我们预期反应 $e^{+}e^{-}\to q\bar{q}$ 的截面趋于 $3\cdot Q^2\cdot 4\pi\alpha^2/3E^2_{\mathrm{cm}}$。习惯上定义
\begin{equation}
  R \equiv \frac{\sigma(e^{+}e^{-}\to\text{强子})}
    {\sigma(e^{+}e^{-}\to\mu^{+}\mu^{-})}
  = \frac{4\pi\alpha^2}{3E_{\mathrm{cm}}^2}
    \quad\text{以 } R \text{ 为单位}
    = \Bigl( \frac{E_{\mathrm{cm}}}{1\ \mathrm{GeV}} \Bigr)^2 \cdot
      \frac{86.8\ \mathrm{nbarns}}{4\pi\alpha^2/3E_{\mathrm{cm}}^2}
      \cdot\ \cdots .
  \label{eq:5.15}
\end{equation}
因此，以 $R$ 为单位的截面值就是其与式~\eqref{eq:5.14} 预言的 $e^{+}e^{-}\to\mu^{+}\mu^{-}$ 截面渐近值的比值。在实验上，最容易测量的量是所有强子的产生总速率。渐近地，我们预期
\begin{equation}
  \frac{\sigma(e^{+}e^{-}\to\text{强子})}{R}
    \to 3\cdot\Bigl( \sum_f Q_f^2 \Bigr),
  \qquad E_{\mathrm{cm}} \to \infty ,
  \label{eq:5.16}
\end{equation}
其中求和遍及所有质量小于 $E_{\mathrm{cm}}/2$ 的夸克。当 $E_{\mathrm{cm}}/2$ 接近某个夸克质量时，强相互作用会引起对该公式的大的偏离。最引人注目的此类效应是束缚态恰好出现在 $E_{\mathrm{cm}} = 2m_q$ 之下，表现为截面中非常尖锐的峰。

图~\ref{fig:5.3} 显示了 2.5 到 40 GeV 之间 $e^{+}e^{-}$ 湮灭为强子截面的实验测量。数据显示了三个不同的区域：一个产生 $u$、$d$ 和 $s$ 夸克对的低能区域；一个高于 $c$ 夸克对产生阈值的区域；以及一个同样高于 $b$ 夸克对阈值的区域。预言 \eqref{eq:5.16} 以一组实线表示；只要能量远离高能近似失效的阈值，它在每个区域都与数据符合得相当好。虚线曲线显示了改进的理论预言，其中包含了来自 QCD 的高阶修正，我们将在第17.2节中讨论。对 $e^{+}e^{-}$ 湮灭截面的这种解释是 QCD 的一个显著成功。特别是，对 \eqref{eq:5.16} 中因子 3 的实验验证是颜色存在的一个证据。

\begin{figure}[t]
  \centering
  \includegraphics[width=0.7\textwidth]{fig5.3.png}
  \caption{反应 $e^{+}e^{-}\to$ 强子的总截面实验测量，数据来自 M.~Swartz 的数据汇编，\emph{Phys.\ Rev.\ D} \textbf{53}，5268 (1996)。各种实验的完整参考文献见该文。如正文所述，测量结果与量子色动力学的理论预言进行了比较。实线是简单预言 \eqref{eq:5.16}。}
  \label{fig:5.3}
\end{figure}

微分截面的角依赖关系也在实验中被观测到。%
\footnote{高能 $e^{+}e^{-}$ 湮灭中强子产生的基本特征由 P.~Duinker 综述，\emph{Rev.\ Mod.\ Phys.} \textbf{54}，325 (1982)。}  高能下，强子以喷注的形式出现，即若干大致沿同一方向运动的强子聚集成的团簇。在大多数情况下有两个动量背靠背的喷注，它们确实具有 $(1 + \cos^2\theta)$ 的角依赖关系。

\section{$e^{+}e^{-}\to\mu^{+}\mu^{-}$：螺旋度结构}

一个反应的未极化截面通常容易计算（也容易测量），但难以理解。$(1 + \cos^2\theta)$ 的角依赖关系从何而来？我们可以通过分别计算每一组自旋取向的 $e^{+}e^{-}\to\mu^{+}\mu^{-}$ 截面来回答这个问题。

首先我们必须选取一个极化态基。为了在高能极限下得到简单的答案，最好的选择是沿粒子运动方向对各自旋进行量子化，即使用具有确定螺旋度的态。回忆在无质量极限下，Dirac 粒子的左手和右手螺旋度态属于 Lorentz 群的不同表示。因此我们可能预期它们独立地行为，事实也确实如此。

在本节中，我们将用螺旋度基以两种不同的方式计算极化的 $e^{+}e^{-}\to\mu^{+}\mu^{-}$ 截面：第一种，使用求迹技巧，但附加螺旋度投影算符以投影出所需的左手或右手旋量；第二种，把这些旋量的显式表达式直接代入我们关于振幅 $\mathcal{M}$ 的公式。在本节的整个过程中，我们都在所有费米子实际上无质量的高能极限下工作。（该计算也可以在较低能量下进行，但要困难得多，而且并没有更多的启发性。）%
\footnote{对这一计算的保留所有质量的一般处理，见 F.~M.~Renard，\emph{Basics of Electron-Positron Collisions}（Editions Fronti\`eres，1981）。}

我们计算极化截面的两种方法的出发点都是振幅
\begin{equation}
  \mathcal{M} = \frac{e^2}{q^2}\
    \bigl( \bar{v}(p')\gamma^\mu u(p) \bigr)\
    \bigl( \bar{u}(k)\gamma_\mu v(k') \bigr) .
  \label{eq:5.19}
\end{equation}
我们使用螺旋度投影算符
\begin{equation}
  P_R = \frac{1+\gamma^5}{2},\qquad P_L = \frac{1-\gamma^5}{2} .
\end{equation}
左手和右手旋量由下式定义
\begin{equation}
  u_L = P_L u, \qquad u_R = P_R u, \qquad v_L = P_L v, \qquad v_R = P_R v .
\end{equation}
为了计算确定螺旋度的振幅，我们在每个旋量前插入投影算符并使用求迹技巧。在对（现在是确定的）螺旋度态求和之后，结果给出极化截面公式。例如，对于过程 $e^{-}_R e^{+}_L \to \mu^{-}_R \mu^{+}_L$，人们得到
\begin{equation}
  \Bigl( \frac{d\sigma}{d\Omega} \Bigr)(e^{-}_R e^{+}_L \to \mu^{-}_R
    \mu^{+}_L)
    = \frac{\alpha^2}{4E_{\mathrm{cm}}^2}\ (1 + \cos\theta)^2 .
  \label{eq:5.21}
\end{equation}
注意截面 \eqref{eq:5.21} 在 $\theta = 180^{\circ}$ 处为零。这正是我们所预期的，因为对于 $\theta = 180^{\circ}$，末态的总角动量与初态的总角动量相反（见图~\ref{fig:5.4}）。

\begin{figure}[t]
  \centering
  \includegraphics[width=0.5\textwidth]{fig5.4.png}
  \caption{角动量守恒要求：如果测量角动量的 $z$ 分量，它必须与初态的取值相同。}
  \label{fig:5.4}
\end{figure}

这完成了我们对极化 $e^{+}e^{-}\to\mu^{+}\mu^{-}$ 截面的第一次计算。我们现在用一种更直接、更有启发性、难度也相当的方式重新进行该计算。我们将直接计算振幅 $\mathcal{M}$（而不是平方振幅），使用旋量和 $\gamma$ 矩阵的显式数值。这种方法确实有缺点：它迫使我们更早地特化到某个特定的参考系，因此失去了显式的 Lorentz 不变性。更实际地说，除了非相对论和极端相对论极限外，它都非常繁琐。

再次考虑振幅
\begin{equation}
  \mathcal{M} = \frac{e^2}{q^2}\
    \bigl( \bar{v}(p')\gamma^\mu u(p) \bigr)\
    \bigl( \bar{u}(k)\gamma_\mu v(k') \bigr) .
  \label{eq:5.25}
\end{equation}
在高能极限下，我们关于 Dirac 旋量的一般表达式变为
\begin{equation}
  u(p) = \begin{pmatrix}
    \sqrt{p\cdot\sigma}\,\xi\\
    \sqrt{p\cdot\bar{\sigma}}\,\xi
  \end{pmatrix}
  \xrightarrow{E\to\infty}
  \sqrt{E}\begin{pmatrix}
    \frac{1+\hat{\mathbf{p}}\cdot\bm{\sigma}}{2}\xi\\
    \frac{1+\hat{\mathbf{p}}\cdot\bm{\sigma}}{2}\xi
  \end{pmatrix};\qquad
  v(p) = \begin{pmatrix}
    \sqrt{p\cdot\sigma}\,\xi\\
    -\sqrt{p\cdot\bar{\sigma}}\,\xi
  \end{pmatrix}
  \xrightarrow{E\to\infty}
  \sqrt{E}\begin{pmatrix}
    \frac{1-\hat{\mathbf{p}}\cdot\bm{\sigma}}{2}\xi\\
    -\frac{1-\hat{\mathbf{p}}\cdot\bm{\sigma}}{2}\xi
  \end{pmatrix} .
  \label{eq:5.26}
\end{equation}
右手旋量满足 $(\hat{\mathbf{p}}\cdot\bm{\sigma})\xi = +\xi$，而左手旋量满足 $(\hat{\mathbf{p}}\cdot\bm{\sigma})\xi = -\xi$。（再次记住，对于反粒子，旋量的手征性与粒子的手征性相反。）我们必须求 $\bar{v}\gamma^\mu u$ 形式的表达式，因此需要恒等式
\begin{equation}
  P_R\gamma^\mu P_L = \gamma^\mu P_L,\qquad
  P_L\gamma^\mu P_R = \gamma^\mu P_R .
\end{equation}
因此我们显式地看到，当两个旋量中一个是左手而另一个是右手时，振幅为零。用第一章的语言来说，把矢量光子耦合到这种旋量乘积的 Clebsch--Gordan 系数为零；这些系数正是矩阵 $\gamma^0\gamma^\mu$ 的分块对角之外的元素（在手征表象中）。

让我们选取 $p$ 和 $p'$ 沿 $\hat{z}$ 方向，并首先考虑电子为右手、正电子为左手的情形。因此对于电子我们有 $\xi = \binom{1}{0}$，对应于 $\hat{z}$ 方向上的自旋向上；而对于正电子我们有 $\xi = \binom{1}{0}$，也对应于（物理上的）$\hat{z}$ 方向自旋向上。两个粒子都满足 $(\hat{\mathbf{p}}\cdot\bm{\sigma})\xi = +\xi$，所以旋量为
\begin{equation}
  u(p) = \sqrt{2E}\begin{pmatrix}\binom{1}{0}\\ \binom{1}{0}\end{pmatrix},
  \qquad v(p') = \sqrt{2E}\begin{pmatrix}\binom{0}{1}\\ -\binom{0}{1}\end{pmatrix}.
  \label{eq:5.28}
\end{equation}
因此矩阵元的电子半边为
\begin{equation}
  \bar{v}(p')\gamma^\mu u(p) = 2E(0, -i)\cdots = -2E(0, 1, i, 0) .
  \label{eq:5.29}
\end{equation}
我们可以这样解释这个表达式：虚光子在 $+z$ 方向具有圆极化；其极化矢量为 $\epsilon^{+} = (1/\sqrt{2})(\hat{x}+i\hat{y})$。

接下来我们必须计算矩阵元的缪子半边。设 $\mu^{-}$ 以与 $\hat{z}$ 轴成角度 $\theta$ 出射，首先考虑它是右手的情形（因此 $\mu^{+}$ 是左手）。为了计算 $\bar{u}(k)\gamma^\mu v(k')$，我们可以回到表达式 \eqref{eq:5.26}，但那样就必须找到与沿缪子动量极化相对应的正确旋量 $\xi$。使用一个技巧要容易得多：由于任何 $\bar{u}\gamma^\mu v$ 形式的表达式都像四矢量一样变换，我们只需旋转结果 \eqref{eq:5.29}。将该矢量在 $xz$ 平面内旋转角度 $\theta$，我们得到
\begin{equation}
  \bar{u}(k)\gamma^\mu v(k') = \bigl[ \bar{v}(k')\gamma^\mu u(k) \bigr]^{*}
    = \bigl[ -2E(0, \cos\theta, i, -\sin\theta) \bigr]^{*}
    = -2E(0, \cos\theta, -i, -\sin\theta) .
  \label{eq:5.30}
\end{equation}
这个矢量也可以解释为虚光子的极化；当它与 \eqref{eq:5.29} 有非零重叠时，我们得到非零振幅。将 \eqref{eq:5.29} 和 \eqref{eq:5.30} 代入 \eqref{eq:5.25}，我们看到振幅为
\begin{equation}
  \mathcal{M}(e^{-}_R e^{+}_L \to \mu^{-}_R\mu^{+}_L)
    = \frac{e^2}{q^2}(2E)^2(-\cos\theta - 1) = -e^2(1+\cos\theta),
  \label{eq:5.31}
\end{equation}
这与 \eqref{eq:1.6} 一致，也与 \eqref{eq:5.21} 一致。现在可以用与上面相同的方式获得这组螺旋度的微分截面，即得到 \eqref{eq:5.21}。

我们可以用类似的方法计算其余三个非零的螺旋度振幅。对于左手电子和右手正电子，我们很容易得到
\begin{equation}
  \bar{v}(p')\gamma^\mu u(p) = -2E(0, 1, -i, 0) = -2E\sqrt{2}\,\epsilon^{\mu}_{-} .
\end{equation}
进行旋转以得到对应于左手 $\mu^{-}$ 和右手 $\mu^{+}$ 的矢量：
\begin{equation}
  \bar{u}(k)\gamma^\mu v(k') = -2E(0, \cos\theta, i, -\sin\theta) .
\end{equation}
以各种方式组合这些碎片得到其余的振幅，
\begin{equation}
  \begin{aligned}
    \mathcal{M}(e^{-}_L e^{+}_R \to \mu^{-}_L\mu^{+}_R) &= -e^2(1+\cos\theta);\\
    \mathcal{M}(e^{-}_L e^{+}_R \to \mu^{-}_R\mu^{+}_L) &=
    \mathcal{M}(e^{-}_R e^{+}_L \to \mu^{-}_L\mu^{+}_R) = -e^2(1-\cos\theta).
  \end{aligned}
  \label{eq:5.32}
\end{equation}

\section{$e^{+}e^{-}\to\mu^{+}\mu^{-}$：非相对论极限}

现在让我们转到能谱的另一端，讨论极端非相对论极限下的反应 $e^{+}e^{-}\to\mu^{+}\mu^{-}$。当 $E$ 仅仅略大于 $m_\mu$ 时，我们之前关于未极化微分截面的结果 \eqref{eq:5.12} 变为
\begin{equation}
  \frac{d\sigma}{d\Omega} \xrightarrow{|\mathbf{k}|\to 0}
    \frac{\alpha^2}{4E^2}\ \frac{|\mathbf{k}|}{E}\ \frac{1}{1}
    = \frac{\alpha^2|\mathbf{k}|}{4E^3} .
  \label{eq:5.33}
\end{equation}
通过用显式旋量求振幅，我们可以重新得到这个结果，并且也能了解反应的自旋依赖关系。

我们再次从矩阵元开始
\begin{equation}
  \mathcal{M} = \frac{e^2}{q^2}\ \bar{v}(p')\gamma^\mu u(p)\
    \bar{u}(k)\gamma_\mu v(k') .
\end{equation}
电子和正电子仍然非常相对论性，因此如果我们选取它们具有确定的螺旋度，这个表达式将最简单。设电子为右手，沿 $+\hat{z}$ 方向运动，正电子为左手，沿 $-\hat{z}$ 方向运动。于是由式~\eqref{eq:5.29} 我们有
\begin{equation}
  \bar{v}(p')\gamma^\mu u(p) = -2E(0, 1, i, 0) .
  \label{eq:5.34}
\end{equation}
在矩阵元的另一半，我们应该使用非相对论表达式
\begin{equation}
  u(k) = \sqrt{m}\begin{pmatrix}\xi\\ \xi\end{pmatrix},\qquad
  v(k') = \sqrt{m}\begin{pmatrix}\xi'\\ -\xi'\end{pmatrix} .
  \label{eq:5.35}
\end{equation}
在本节的讨论中请记住，旋量 $\xi'$ 给出反粒子的翻转自旋。暂时让缪子旋量 $\xi$ 和 $\xi'$ 不定，我们可以很容易地计算出
\begin{equation}
  \bar{u}(k)\gamma^\mu v(k') = \cdots = -m\bigl( \xi^{\dagger}\xi',
    -\xi^{\dagger}\bm{\sigma}\xi' \bigr) .
  \label{eq:5.36}
\end{equation}
为了求 $\mathcal{M}$，我们只需将 \eqref{eq:5.34} 与 \eqref{eq:5.36} 点乘并乘以 $e^2/q^2 = e^2/4m^2$。结果是
\begin{equation}
  \mathcal{M}(e^{-}_R e^{+}_L \to \mu^{+}\mu^{-})
    = -2e^2\,\xi^{\dagger}\begin{pmatrix}0&0\\0&1\end{pmatrix}\xi' .
  \label{eq:5.37}
\end{equation}
由于这个表达式中没有角依赖关系，缪子向任何方向出射的可能性都相同。更精确地说，它们以 $S$ 波发射；其轨道角动量为零。因此角动量守恒要求末态的总自旋等于 1，事实上，除非缪子和反缪子都沿 $\hat{z}$ 轴自旋向上，否则矩阵乘积为零（见图~\ref{fig:5.5}）。

\begin{figure}[t]
  \centering
  \includegraphics[width=0.5\textwidth]{fig5.5.png}
  \caption{在非相对论极限下，系统的总自旋守恒，因此产生的缪子都沿 $\hat{z}$ 轴自旋向上。}
  \label{fig:5.5}
\end{figure}

为了求这个过程的总速率，我们对缪子自旋求和得到 $|\mathcal{M}|^2 = 4e^4$，从而得到截面
\begin{equation}
  \frac{d\sigma}{d\Omega}(e^{-}_R e^{+}_L \to \mu^{+}\mu^{-})
    = \frac{\alpha^2|\mathbf{k}|}{4m_\mu^2 E^2} .
  \label{eq:5.38}
\end{equation}
同样的表达式也适用于左手电子和右手正电子的情形。因此自旋平均的截面正是该表达式的 $2\cdot(1/4)$ 倍，与 \eqref{eq:5.33} 一致。

\subsection{束缚态}

到目前为止，我们都把散射过程的初态和末态视为孤立单粒子的态。然而，非常接近阈值时，缪子之间的 Coulomb 吸引应当成为一个重要的效应。就在阈值之下，我们仍然可以在电磁束缚态中形成 $\mu^{+}\mu^{-}$ 对。

量子场论中束缚态的处理是一个丰富而复杂的课题，但它主要超出了本书的范围。%
\footnote{例如，见 C.~Itzykson 和 J.-B.~Zuber 的综述，\emph{Quantum Field Theory}（McGraw-Hill，1980），以及其中的参考文献。} 幸运的是，自然界中许多熟悉的束缚系统都可以（至少在一级近似下很好地）作为非相对论系统来处理，在这些系统中内部运动是缓慢的。从真空中产生组成粒子的过程仍然是一种相对论效应，需要量子场论才能恰当描述。在本节中，我们将发展一个形式体系，用于计算双粒子非相对论束缚态的产生和湮灭振幅。我们从计算 $e^{+}e^{-}$ 湮灭中产生 $\mu^{+}\mu^{-}$ 束缚态的截面开始。

首先考虑电子和正电子自旋都指向同一方向的情形。一个费米子及其反费米子的一般束缚态具有总自旋 $S = 0$ 或 $S = 1$，以及空间波函数 $\psi(\mathbf{r})$。产生束缚态 $B$ 的振幅的非相对论极限正比于原点处的波函数，
\begin{equation}
  \mathcal{M}(e^{+}e^{-}\to B) \propto \psi(0) .
  \label{eq:5.50}
\end{equation}
完成计算（细节见正文），我们得到
\begin{equation}
  \mathcal{M}(e^{+}e^{-}\to B) = 2e^2\sqrt{2M}\,\psi(0)\,
    \epsilon(\mathbf{n})\cdot\hat{z}
  \label{eq:5.51}
\end{equation}
这是对极化矢量为 $\bm{\epsilon}$ 的自旋为 1 的束缚态的情形，其中 $M = 2m_\mu$ 是束缚态质量，$\mathbf{n}$ 是缪子动量的方向。假设电子束未极化，并对 $\mathbf{n}$ 的三个可能方向求和 \eqref{eq:5.52}，我们得到束缚态产生的总截面的如下表达式：
\begin{equation}
  \sigma(e^{+}e^{-}\to B) =
    \frac{(2\pi)^4}{2E^2\,2E}\ \frac{1}{2}\ \delta(P^0 - M^0)\
    \frac{1}{3}(4e^4)|\psi(0)|^2 .
  \label{eq:5.53}
\end{equation}
注意单体相空间积分只能消去四个 delta 函数中的三个。习惯上用 $\delta(P^0 - K^0) = 2K^0\,\delta(P^2 - K^2)$ 重写最后一个 delta 函数。于是
\begin{equation}
  \sigma(e^{+}e^{-}\to B) = 64\pi^3\alpha^2\ \frac{|\psi(0)|^2}{M}\
    \delta(E^2_{\mathrm{cm}} - M^2) .
  \label{eq:5.54}
\end{equation}
最后一个 delta 函数强制了总质心能量必须等于束缚态质量的约束；因此，束缚态在 $e^{+}e^{-}$ 湮灭中作为共振产生。如果束缚态具有有限寿命，这个 delta 函数将被展宽为共振峰。在实践中，$e^{+}e^{-}$ 束能量的固有展宽往往是更重要的展宽机制。无论哪种情况，\eqref{eq:5.54} 都正确预言了共振峰下的面积。

如果束缚态 $B$ 能从 $e^{+}e^{-}$ 产生，那么它也能湮灭回 $e^{+}e^{-}$，或湮灭为任何其他足够轻的轻子对。根据 \eqref{eq:4.86}，这种衰变方式的总宽度由下式给出
\begin{equation}
  \Gamma(B \to e^{+}e^{-}) = \frac{1}{2M}\ \int d\Pi_2\ |\mathcal{M}|^2 ,
  \label{eq:5.55}
\end{equation}
其中 $\mathcal{M}$ 正是我们用来计算 $B$ 产生的矩阵元 \eqref{eq:5.51} 的复共轭。因此
\begin{equation}
  \Gamma(B \to e^{+}e^{-}) = 16\pi\alpha^2\ \frac{|\psi(0)|^2}{M^2}\
    \cdot\ 3\ \cdot\ \cdots .
  \label{eq:5.57}
\end{equation}
这个公式是对可能的末态电子极化态求和的。最容易的求值方法是平均 $\mathbf{n}$ 的三个可能值。我们由此得到 \eqref{eq:5.57}。$B$ 衰变宽度的公式与产生截面的公式非常相似，这并不奇怪：两个计算都涉及同一个矩阵元的平方，并对初态和末态极化求和。两个计算的不同仅在于我们如何构造极化平均，以及相空间因子。根据这一逻辑，我们发现的两个量之间的关系，
\begin{equation}
  \sigma(e^{+}e^{-}\to B) = 4\pi^2\cdot\frac{3}{M^2}\cdot
    \Gamma(B\to e^{+}e^{-})\ \delta(E^2_{\mathrm{cm}} - M^2),
  \label{eq:5.58}
\end{equation}
是非常普遍的，并且完全独立于矩阵元计算的细节。\eqref{eq:5.58} 中的因子 3 来自对 $\mathbf{n}$ 的方向平均；对于自旋为 $J$ 的束缚态，这个因子将是 $(2J+1)$。

这个形式体系最著名的应用不是针对缪子的束缚态，而是针对夸克的束缚态：夸克偶素。我们曾在图~\ref{fig:5.3} 中看到过 $q\bar{q}$ 束缚态（例如 $J/\psi$ 和 $\Upsilon$）的实验证据。（共振峰太高太窄，无法在图中显示，但它们的尺寸已被仔细测量过。）方程 \eqref{eq:5.54} 和 \eqref{eq:5.57} 必须乘以颜色因子 3，才能给出自旋为 1 的 $q\bar{q}$ 束缚态的产生截面和衰变宽度。$q\bar{q}$ 波函数在原点处的值 $\psi(0)$ 无法从第一性原理计算，但可以用带唯象选取势的 $q\bar{q}$ 谱的非相对论模型来估计。或者，我们可以用公式
\begin{equation}
  \Gamma\bigl( B(q\bar{q}) \to e^{+}e^{-} \bigr) = 16\pi\alpha^2 Q^2
    \frac{|\psi(0)|^2}{M^2}
  \label{eq:5.59}
\end{equation}
来测量 $q\bar{q}$ 束缚态的 $\psi(0)$。例如，$s\bar{s}$ 的 $1S$ 自旋为 1 的态，即 $\phi$ 介子，其 $e^{+}e^{-}$ 分宽度为 1.4 keV，质量为 1.02 GeV。由此我们可以推断 $|\psi(0)|^2 = (1.2\ \mathrm{fm})^{-3}$。这个结果在物理上是合理的，因为强子的尺寸通常 $\sim 1$ fm。

本节的观点与前面几节相当不同：我们不是从第一性原理计算一切，而是用一点量子场论和一点非相对论量子力学拼凑出一个近似公式。然而，原则上我们可以在相对论形式体系中完全处理束缚态。考虑一个 $e^{+}e^{-}$ 对湮灭形成一个 $\mu^{+}\mu^{-}$ 束缚态，随后该束缚态衰变回 $e^{+}e^{-}$。在我们目前的形式体系中，我们可以用下列图形来表示这个过程
\begin{equation}
  e^{+}e^{-} \to B \to e^{+}e^{-} .
\end{equation}
净过程就是简单的 $e^{+}e^{-}\to e^{+}e^{-}$（巴巴散射）。如果我们试图直接在 QED 微扰论中计算巴巴散射截面，会发生什么？显然，在树级图中没有 $\mu^{+}\mu^{-}$ 的贡献。然而，当我们进行到微扰级数的更高阶时，我们（除其他图外）会发现如下一组图，其中两条电子线之间交换一个 $\mu^{+}\mu^{-}$ 圈。在大多数 $E_{\mathrm{cm}}$ 值下，这些图对树级表达式只给出小的修正。但当 $E_{\mathrm{cm}}$ 接近 $\mu^{+}\mu^{-}$ 阈值时，涉及缪子圈内交换光子的图包含缪子之间的 Coulomb 相互作用，因此变得相当大。人们必须对所有这样的图求和，并且可以证明这种求和等价于求解非相对论的 Schr\"odinger 方程。%
\footnote{这一分析在 Berestetskii、Lifshitz 和 Pitaevskii (1982) 中完成。} 最终预言是截面包含一个共振峰，其面积由 \eqref{eq:5.54} 给出，其宽度由 \eqref{eq:5.57} 给出。

\section{交叉对称性}

\subsection{电子-缪子散射}

既然我们已经完成了对过程 $e^{+}e^{-}\to\mu^{+}\mu^{-}$ 的讨论，让我们考虑一个不同但密切相关的 QED 过程：电子-缪子散射，即 $e^{-}\mu^{-}\to e^{-}\mu^{-}$。其最低阶费曼图正是把上一个图侧转过来：
\begin{equation}
  i\mathcal{M} = \frac{(-ie)^2}{q^2}\ \bar{u}(p_1')\gamma^\mu u(p_1)\
    \bar{u}(p_2')\gamma_\mu u(p_2) .
\end{equation}
当我们计算对自旋求平均并求和的平方振幅时，过程 $e^{+}e^{-}\to\mu^{+}\mu^{-}$ 与 $e^{-}\mu^{-}\to e^{-}\mu^{-}$ 之间的关系就变得清楚了：
\begin{equation}
  \frac{1}{4}\sum_{\text{spins}}|\mathcal{M}|^2
    = \frac{e^4}{4q^4}\ \mathrm{tr}\bigl[
      (\slashed{p}_1' + m_e)\gamma^\mu(\slashed{p}_1 + m_e)\gamma^\nu \bigr]\
      \mathrm{tr}\bigl[ (\slashed{p}_2' + m_\mu)\gamma_\mu(\slashed{p}_2
        + m_\mu)\gamma_\nu \bigr] .
\end{equation}
这与我们关于 $e^{+}e^{-}\to\mu^{+}\mu^{-}$ 的结果 \eqref{eq:5.4} 完全相同，只需作替换
\begin{equation}
  p \to p_1,\qquad p' \to -p_1',\qquad k \to p_2',\qquad k' \to -p_2 .
  \label{eq:5.60}
\end{equation}
因此，不必从头求迹，我们只需在我们之前的结果式~\eqref{eq:5.10} 中作同样的替换。令 $m_e = 0$，我们得到
\begin{equation}
  \frac{1}{4}\sum_{\text{spins}}|\mathcal{M}|^2
    = \frac{e^4}{4q^4}\ \Bigl[ (p_1\cdot p_2)(p_1'\cdot p_2')
      + (p_1\cdot p_2')(p_1'\cdot p_2) - m_\mu^2(p_1\cdot p_1') \Bigr] .
  \label{eq:5.61}
\end{equation}
为了求这个表达式，我们必须确定运动学，它将完全不同。在质心系中工作，我们作如下赋值：
\begin{equation}
  \begin{aligned}
    p_1 &= (k, \mathbf{k}),\qquad p_1' = (k, \mathbf{k}');\qquad
    p_2 = (E, -\mathbf{k}),\qquad p_2' = (E, -\mathbf{k}').
  \end{aligned}
\end{equation}
我们所需的组合是
\begin{equation}
  \begin{aligned}
    p_1\cdot p_2 = p_1'\cdot p_2' &= k(E+k);\\
    p_1\cdot p_1' &= k^2(1-\cos\theta);\\
    p_1\cdot p_2' = p_1'\cdot p_2 &= k(E+k\cos\theta);\\
    q^2 &= -2p_1\cdot p_1' = -2k^2(1-\cos\theta) .
  \end{aligned}
\end{equation}
我们关于平方矩阵元的表达式现在变为
\begin{equation}
  \frac{1}{4}\sum_{\text{spins}}|\mathcal{M}|^2
    = \frac{e^4}{4k^2E^2(1-\cos\theta)^2}\ \Bigl[ k^2(E+k)^2
      + k^2(E+k\cos\theta)^2 - m_\mu^2 k^2(1-\cos\theta) \Bigr] .
  \label{eq:5.62}
\end{equation}
为了从这个表达式求截面，我们使用式~\eqref{eq:4.84}，在一个粒子无质量的情形下它取简单的形式
\begin{equation}
  \Bigl( \frac{d\sigma}{d\Omega} \Bigr)_{\mathrm{CM}}
    = \frac{|\mathcal{M}|^2}{64\pi^2(E+k)^2} .
  \label{eq:5.63}
\end{equation}
因此我们得到了质心系中未极化电子-缪子散射的结果：
\begin{equation}
  \Bigl( \frac{d\sigma}{d\Omega} \Bigr)_{\mathrm{CM}}
    = \frac{e^4}{64\pi^2(E+k)^2}\ \frac{1}{4k^2E^2(1-\cos\theta)^2}\
      \Bigl[ \cdots \Bigr] ,
  \label{eq:5.64}
\end{equation}
其中 $k = \sqrt{E^2-m_\mu^2}$。在可以令 $m_\mu = 0$ 的高能极限下，微分截面变为
\begin{equation}
  \frac{d\sigma}{d\Omega} \xrightarrow{E\gg m_\mu}
    \frac{\alpha^2}{4E_{\mathrm{cm}}^2}\ \frac{1}{\sin^4(\theta/2)}
      \Bigl( \frac{1+\cos^2\theta}{4} \Bigr) .
  \label{eq:5.65}
\end{equation}

\subsection{交叉对称性与曼德尔施塔姆变量}

上一小节的操作是散射振幅之间一个更一般关系的特例，称为\emph{交叉对称性}。考虑一个一般的两体散射过程。习惯上引入曼德尔施塔姆变量
\begin{equation}
  s = (p+p')^2,\qquad t = (k-p)^2,\qquad u = (k'-p)^2 ,
  \label{eq:5.69}
\end{equation}
其中 $p$ 和 $p'$ 是初态动量，$k$ 和 $k'$ 是末态动量。这三个变量满足
\begin{equation}
  s + t + u = \sum_{i=1}^{4} m_i^2 ,
  \label{eq:5.73}
\end{equation}
其中求和遍及四个外线粒子。这个恒等式很容易证明：把式~\eqref{eq:5.69} 右边的各项加起来，并应用动量守恒 $(p+p'-k-k')^2 = 0$。

交叉对称性把一个过程的振幅与相关过程的振幅联系起来，在相关过程中一些初态粒子被移到末态（把它们的动量替换为其相反数）。例如，用曼德尔施塔姆变量表示的 $e^{+}e^{-}\to\mu^{+}\mu^{-}$ 振幅可以通过替换从 $e^{-}\mu^{-}\to e^{-}\mu^{-}$ 的振幅得到
\begin{equation}
  p \leftrightarrow k,\qquad p' \to -p',\qquad k' \to -k' .
  \label{eq:5.68}
\end{equation}
在四个质量都等于 $m$ 的质心系中，我们有
\begin{equation}
  s = (p+p')^2 = (2E)^2 = E^2_{\mathrm{cm}},\qquad
  t = (k-p)^2 = -2p^2(1-\cos\theta),\qquad
  u = (k'-p)^2 = -2p^2(1+\cos\theta) .
  \label{eq:5.72}
\end{equation}
因此我们看到，当 $\theta\to 0$ 时 $t\to 0$，而当 $\theta\to\pi$ 时 $u\to 0$。（当质量不全相等时，$t$ 和 $u$ 的极限值会有少许移动。）注意由 \eqref{eq:5.72}，当四个粒子质量都为 $m$ 时，曼德尔施塔姆变量之和为 $s+t+u = 4E^2 - 4p^2 = 4m^2$。这是更一般的式 \eqref{eq:5.73} 的一个特例。

\section{康普顿散射}

我们现在转而考虑一个稍有不同 QED 过程：康普顿散射，即 $e^{-}\gamma\to e^{-}\gamma$。我们将计算该反应对 $\alpha$ 最低阶的未极化截面。该计算将用到我们迄今发展的所有工具，包括上一节的曼德尔施塔姆变量。我们还将发展一些处理外线光子的新技术。

这是我们第一个涉及两个费曼图的计算例子：
\begin{equation}
  i\mathcal{M} = \bar{u}(p')(-ie\gamma^\mu)\epsilon^{*}_{\mu}(k')\
    \frac{i(\slashed{p}+\slashed{k}+m)}{(p+k)^2-m^2}\
    (-ie\gamma^\nu)\epsilon_\nu(k)u(p)
    + \bar{u}(p')(-ie\gamma^\mu)\epsilon_\nu(k)\
      \frac{i(\slashed{p}-\slashed{k}'+m)}{(p-k')^2-m^2}\
      (-ie\gamma^\nu)\epsilon^{*}_\mu(k')u(p) .
\end{equation}
如常，费曼规则精确地告诉我们如何写出 $\mathcal{M}$ 的表达式。注意由于两个图的费米子部分是相同的，两项之间没有相对的负号。用 $\epsilon_\mu(k)$ 和 $\epsilon^{*}_\mu(k')$ 表示初态和末态光子的极化矢量，我们有
\begin{equation}
  \begin{aligned}
    i\mathcal{M} &= -ie^2\ \epsilon^{*}_{\nu}(k')\epsilon_\mu(k)\
      \bar{u}(p')\ \Bigl[
        \frac{\gamma^\mu(\slashed{p}+\slashed{k}+m)\gamma^\nu}{(p+k)^2-m^2}
        + \frac{\gamma^\nu(\slashed{p}-\slashed{k}'+m)\gamma^\mu}
          {(p-k')^2-m^2} \Bigr]\ u(p) .
  \end{aligned}
  \label{eq:5.74}
\end{equation}
在平方这个表达式之前，我们可以作一些简化。由于 $p^2 = m^2$ 和 $k^2 = 0$，传播子的分母为
\begin{equation}
  (p+k)^2 - m^2 = 2p\cdot k
  \qquad\text{和}\qquad
  (p-k')^2 - m^2 = -2p'\cdot k' .
\end{equation}
为了简化分子，我们使用一点 Dirac 代数：
\begin{equation}
  (\slashed{p}+m)\gamma^\nu u(p) = (2p^\nu - \gamma^\nu\slashed{p}
    + \gamma^\nu m)u(p) = 2p^\nu u(p) .
\end{equation}
对每个传播子的分子使用这个技巧，我们得到
\begin{equation}
  i\mathcal{M} = -ie^2\ \epsilon^{*}_{\nu}(k')\epsilon_\mu(k)\
    \bar{u}(p')\ \Bigl[ \frac{\gamma^\mu\slashed{k}\gamma^\nu}{2p\cdot k}
      - \frac{\gamma^\nu\slashed{k}'\gamma^\mu}{2p'\cdot k'} \Bigr]\ u(p) .
  \label{eq:5.74b}
\end{equation}

\subsection{光子极化求和}

计算的下一步将是对 $\mathcal{M}$ 的这个表达式平方，并对电子和光子的极化态求和（或求平均）。对电子极化的求和可以像以前一样用恒等式 $\sum u(p)\bar{u}(p) = \slashed{p} + m$ 来完成。幸运的是，对光子极化矢量求和也有类似的技巧。正确的做法是作替换
\begin{equation}
  \sum_{\text{polarizations}}\ \epsilon^{*}_\mu(k)\epsilon_\nu(k)
    \to -g_{\mu\nu} .
  \label{eq:5.75}
\end{equation}
箭头表示这并不是一个真正的等式。尽管如此，只要两边都点乘到 QED 振幅 $\mathcal{M}$ 表达式的其余部分，这个替换就是有效的。

为了推导这个公式，让我们考虑一个涉及动量为 $k$ 的外线光子的任意 QED 过程：
\begin{equation}
  \mathcal{M}(k) = i\mathcal{M}^{\mu}(k)\epsilon^{*}_\mu(k) .
  \label{eq:5.76}
\end{equation}
由于振幅总是包含 $\epsilon^{*}_\mu(k)$，我们提取出这个因子，并把振幅 $\mathcal{M}$ 的其余部分定义为 $\mathcal{M}^{\mu}(k)$。截面将正比于
\begin{equation}
  \sum_{\epsilon} \bigl| \epsilon^{*}_\mu(k)\mathcal{M}^{\mu}(k) \bigr|^2
    = \sum_{\epsilon} \epsilon^{*}_\mu(k)\epsilon_\nu(k)\,
      \mathcal{M}^{\mu}(k)\mathcal{M}^{\nu*}(k) .
\end{equation}
为简单起见，我们把 $k$ 取向为 3 方向：$k^\mu = (k, 0, 0, k)$。于是我们要对其求和的这两个横极化矢量可以选为
\begin{equation}
  \epsilon_1 = (0, 1, 0, 0);\qquad \epsilon_2 = (0, 0, 1, 0) .
\end{equation}
在这些约定下，我们有
\begin{equation}
  \sum_{\epsilon} \bigl| \epsilon^{*}_\mu(k)\mathcal{M}^{\mu}(k) \bigr|^2
    = |\mathcal{M}^1(k)|^2 + |\mathcal{M}^2(k)|^2 .
  \label{eq:5.77}
\end{equation}
现在回忆第四章，外线光子是由相互作用项 $\int d^4x\ e j^\mu A_\mu$ 产生的，其中 $j^\mu = \bar{\psi}\gamma^\mu\psi$ 是 Dirac 矢量流。因此我们预期 $\mathcal{M}^{\mu}(k)$ 由 Heisenberg 场 $j^\mu$ 的矩阵元给出：
\begin{equation}
  \mathcal{M}^{\mu}(k) = \int d^4x\ e^{ik\cdot x}\ \langle f|j^\mu(x)|i\rangle ,
  \label{eq:5.78}
\end{equation}
其中初态和末态包含除所讨论的光子外的所有粒子。

从经典运动方程，我们知道流 $j^\mu$ 是守恒的：$\partial_\mu j^\mu(x) = 0$。只要这个性质在量子理论中仍然成立，我们就可以把 $k_\mu$ 点乘到 \eqref{eq:5.78} 上得到
\begin{equation}
  k_\mu\mathcal{M}^{\mu}(k) = 0 .
  \label{eq:5.79}
\end{equation}
当极化矢量 $\epsilon^\mu(k)$ 被替换为 $k^\mu$ 时，振幅 $\mathcal{M}$ 为零。这个著名的关系被称为\emph{沃德恒等式}。它本质上是流守恒的表述，而流守恒是 QED 规范对称性 \eqref{eq:4.6} 的推论。我们将在第7.4节给出沃德恒等式的正式证明，并讨论这个快速``推导''中一带而过的若干微妙之处。

显式地验证 \eqref{eq:5.74b} 中给出的康普顿振幅遵从沃德恒等式是有用的。为此，把 $\epsilon_\mu(k)$ 替换为 $k_\mu$，或把 $\epsilon^{*}_\mu(k')$ 替换为 $k'_\mu$，并处理 Dirac 矩阵乘积。无论哪种情况（经过一点代数运算），来自两个图的项相互抵消而为零。

回到我们对极化求和公式 \eqref{eq:5.75} 的推导，我们注意到对于 $k^\mu = (k, 0, 0, k)$，沃德恒等式取如下形式
\begin{equation}
  k\mathcal{M}^0(k) - k\mathcal{M}^3(k) = 0 .
  \label{eq:5.80}
\end{equation}
因此 $\mathcal{M}^0 = \mathcal{M}^3$，于是我们有
\begin{equation}
  \begin{aligned}
    \sum_{\epsilon} \epsilon^{*}_\mu(k)\epsilon_\nu(k)\,
      \mathcal{M}^{\mu}(k)\mathcal{M}^{\nu*}(k)
    &= |\mathcal{M}^1|^2 + |\mathcal{M}^2|^2\\
    &= |\mathcal{M}^1|^2 + |\mathcal{M}^2|^2 + |\mathcal{M}^3|^2
      - |\mathcal{M}^0|^2\\
    &= -g_{\mu\nu}\mathcal{M}^{\mu}(k)\mathcal{M}^{\nu*}(k) .
  \end{aligned}
  \label{eq:5.81}
\end{equation}
也就是说，我们可以通过把 $\sum \epsilon^{*}_\mu\epsilon_\nu$ 替换为 $-g_{\mu\nu}$ 来对外线光子极化求和。

注意这证明了（有待我们对沃德恒等式的一般证明）非物理的类时光子和纵光子可以在 QED 计算中被一致地忽略，因为无论如何产生这些态的平方振幅都相互抵消从而给出零总概率。类时光子态的负模方——这个在式~\eqref{eq:4.132} 之后的讨论中困扰过我们的性质——在这种抵消中起着至关重要的作用。

\subsection{克莱因-仁科公式}

康普顿散射截面计算的其余部分很直接，不过有条理一些会有所帮助。我们希望将平方振幅对初态电子和光子极化求平均，并对末态电子和光子极化求和。从 $\mathcal{M}$ 的表达式 \eqref{eq:5.74b} 出发，我们得到
\begin{equation}
  \begin{aligned}
    \frac{1}{4}\sum_{\text{spins}}|\mathcal{M}|^2
    &= \frac{e^4}{4}\ g_{\mu\nu}g_{\rho\sigma}\ \mathrm{tr}\Bigl\{
      (\slashed{p}'+m)\ \Bigl[
        \frac{\gamma^\mu\slashed{k}\gamma^\nu}{2p\cdot k}
        - \frac{\gamma^\nu\slashed{k}'\gamma^\mu}{2p'\cdot k'} \Bigr]\
      (\slashed{p}+m)\ \Bigl[
        \frac{\gamma^\rho\slashed{k}\gamma^\sigma}{2p\cdot k}
        - \frac{\gamma^\sigma\slashed{k}'\gamma^\rho}{2p'\cdot k'} \Bigr]
      \Bigr\} .
  \end{aligned}
  \label{eq:5.82}
\end{equation}
该迹分成四项，I、II、III 和 IV：
\begin{equation}
  \frac{1}{4}\sum_{\text{spins}}|\mathcal{M}|^2
    = \frac{e^4}{4}\ \Bigl[ \frac{I}{(2p\cdot k)^2}
      + \frac{II + III}{(2p\cdot k)(2p'\cdot k')}
      + \frac{IV}{(2p'\cdot k')^2} \Bigr],
  \label{eq:5.83}
\end{equation}
其中 I、II、III 和 IV 是复杂的迹。注意如果把 $k$ 替换为 $-k'$，IV 与 I 相同。此外，由于我们可以在迹内颠倒 $\gamma$ 矩阵的顺序（式~\eqref{eq:5.7}），我们看到 $II = III$。因此我们只需计算 I 和 II。

这些迹中的第一个是
\begin{equation}
  I = \mathrm{tr}\bigl[ (\slashed{p}'+m)(\gamma^\mu\slashed{k}\gamma^\nu
    + 2\gamma^\nu p^\nu)(\slashed{p}+m)(\gamma_\mu\slashed{k}\gamma_\nu
      + 2\gamma_\nu p_\nu) \bigr] .
  \label{eq:5.84}
\end{equation}
迹内有 16 项，但其中一半含有奇数个 $\gamma$ 矩阵因而为零。我们现在必须逐一求其余八项。例如，
\begin{equation}
  \mathrm{tr}[\slashed{p}'\gamma^\mu\slashed{k}\gamma^\nu\slashed{p}
    \gamma_\mu\slashed{k}\gamma_\nu]
  = \mathrm{tr}\bigl[ (-2\slashed{k})\gamma^\nu(-2\slashed{k}')\gamma_\nu
    \bigr]
  = \mathrm{tr}\bigl[ 4\slashed{k}\slashed{k}'(2p\cdot k'-\cdots) \bigr],
\end{equation}
等等。完成所有的迹并合并各项，我们得到未极化康普顿截面的标准结果，
\begin{equation}
  \Bigl( \frac{d\sigma}{d\cos\theta} \Bigr)_{\mathrm{lab}}
    = \frac{\pi\alpha^2}{m^2}\ \Bigl(
      \frac{\omega'}{\omega} \Bigr)^2\ \Bigl[
        \frac{\omega'}{\omega} + \frac{\omega}{\omega'}
        - \sin^2\theta \Bigr] ,
  \label{eq:5.94}
\end{equation}
其中 $\omega$ 和 $\omega'$ 是实验室系中初态和末态光子的能量，它们之间的关系为
\begin{equation}
  \frac{1}{\omega'} - \frac{1}{\omega} = \frac{1}{m}(1-\cos\theta) .
  \label{eq:5.95}
\end{equation}
这就是克莱因-仁科公式。

微分截面在 $\theta\to\pi$ 处的奇异性值得一些注意。奇异项来自 $t$ 道图的平方，其分母在后向角处变小。让我们试着理解这个奇异性的物理。奇异项来自 $t$ 道图的平方，
\begin{equation}
  (e^{-}(p)\ \text{后向发射末态光子}).
\end{equation}
振幅在 $\theta\approx\pi$ 处很大，因为此时传播子的分母（$\sim m^2$）与 $s$ 相比很小。更精确地说，定义 $\chi$ 为
\begin{equation}
  \chi \equiv \pi - \theta .
\end{equation}
我们感兴趣的是略大于 $m/\omega$ 但仍然足够小以至于可以近似 $1-\cos\chi \approx \chi^2/2$ 的 $\chi$ 值。对于这个范围内的 $\chi$，分母为
\begin{equation}
  (p-k')^2 - m^2 = -2p\cdot k' \approx -2\omega^2\Bigl( \frac{m^2}{\omega^2}
    + 1 - \cos\chi \Bigr) \approx -(\omega^2\chi^2 + m^2) .
  \label{eq:5.98}
\end{equation}
在 $\chi$ 的很大取值范围内，这与 $s$ 相比很小，因此总截面得到增强。

回顾 \eqref{eq:5.93}，我们看到对于满足 $m/\omega \ll \chi \ll 1$ 的 $\chi$，平方振幅正比于 $1/\chi^2$，因此我们预期 $\mathcal{M} \propto 1/\chi$。但我们已经看到 $\mathcal{M}$ 的分母正比于 $\chi^2$，所以分子中必定有一个补偿的 $\chi$ 因子。通过考察一组特定的电子和光子极化的振幅，我们可以理解这个因子的物理起源。

假设初态电子是右手的。\eqref{eq:5.97} 的主要项来自传播子分子中涉及 $(\slashed{p}-\slashed{k}')$ 的项。由于在 \eqref{eq:5.97} 中这一项在 $\bar{u}$ 和 $u$ 之间包含三个 $\gamma$ 矩阵，末态电子也必定是右手的。因此振幅为
\begin{equation}
  i\mathcal{M} = -ie^2\ \epsilon^{*}_{\mu}(k')\epsilon_\nu(k)\
    \bar{u}_R(p')[\cdots]u_R(p) ,
  \label{eq:5.99}
\end{equation}
其中被省略的因子是极化矢量和传播子的分子。如果初态光子是左手的，具有 $\epsilon^\mu = (1/\sqrt{2})(0, 1, -i, 0)$，那么组合 $\bar{u}_R(p')\slashed{\epsilon}(k)u_R(p)$ 为零。因此初态光子必定是右手的。类似地，除非末态光子是右手的，否则振幅为零。这组极化的运动学情形如图~\ref{fig:5.6} 所示。

\begin{figure}[t]
  \centering
  \includegraphics[width=0.5\textwidth]{fig5.6.png}
  \caption{在高能极限下，末态光子最可能在后向角出射。由于螺旋度守恒，一个单位的自旋角动量转化为轨道角动量。}
  \label{fig:5.6}
\end{figure}

继续我们的计算，让我们考虑 \eqref{eq:5.99} 中传播子的分子。对于感兴趣范围内的 $\chi$，主要项为
\begin{equation}
  \bm{\sigma}\cdot(\mathbf{p}-\mathbf{k}') \approx \bm{\sigma}\cdot\hat{x}\,
    \omega\chi .
\end{equation}
这就是上面预期的 $\chi$ 因子。它表明末态是 $p$ 波，这是角动量守恒所要求的。把所有碎片组装起来，我们得到
\begin{equation}
  \mathcal{M}(e^{-}_R\gamma_R \to e^{-}_R\gamma_R)
    \approx e^2\sqrt{2E}\ \sqrt{2E'}\ \frac{\omega\chi}{\omega^2\chi^2+m^2}\
    \cdots .
  \label{eq:5.101}
\end{equation}
在所有初态和末态粒子都是左手的情形下，我们会得到相同的结果。

注意对于正后向散射，$\chi = 0$，矩阵元 \eqref{eq:5.101} 由于分子中的角动量为零而为零。因此，在非常接近后向的角度，我们还应该考虑 \eqref{eq:5.97} 中传播子分子里的质量项。这一项只包含两个伽马矩阵，因此把右手电子转化为左手电子。通过与导出式~\eqref{eq:5.101} 的分析类似的分析，我们可以看出这个振幅只有在初态光子是左手而末态光子是右手时才非零。更详细地跟进这一分析，我们得到
\begin{equation}
  \mathcal{M}(e^{-}_R\gamma_L \to e^{-}_L\gamma_R) = \frac{4e^2 m/\omega}
    {\chi^2 + m^2/\omega^2}
  \label{eq:5.102}
\end{equation}
所有四个螺旋度都反转的反应给出相同的矩阵元。

为了将这一结果与我们之前的计算进行比较，我们应该把 \eqref{eq:5.101} 和 \eqref{eq:5.102} 对截面的贡献以及涉及初态左手电子的反应的相等贡献相加，并除以 4 以对初态自旋求平均。于是未极化微分截面应为
\begin{equation}
  \begin{aligned}
    \frac{d\sigma}{d\cos\theta} &= \frac{1}{4}\cdot\frac{1}{2E\,2\omega}\,
      \frac{1}{(2\pi)^4}\ \frac{8e^4\chi^2}{(E+\omega)(\chi^2+m^2/\omega^2)^2}
      + \frac{8e^4 m^2/\omega^2}{(\chi^2+m^2/\omega^2)^2}\\
    &= \frac{4\pi\alpha^2}{s}\ \frac{1}{\chi^2+4m^2/s}\ \cdots ,
  \end{aligned}
  \label{eq:5.103}
\end{equation}
这与式~\eqref{eq:5.94} 精确一致。

螺旋度翻转过程 \eqref{eq:5.102} 恰在运动学端点处的重要性有一个有趣的实验推论。考虑\emph{逆康普顿散射}过程：一束高能电子与一束低能光子（例如激光束）碰撞产生一束高能光子。设电子的能量为 $E$，激光光子的能量为 $\omega$，设散射光子的能量为 $E' = yE$，并为了简单起见假设 $s = 4E\omega \gg m^2$。那么我们刚才所做的计算适用于这种情况，其中最高能量的光子来自在质心系中恰好后向的散射。通过在质心系和实验室系中计算 $2k\cdot k'$，容易证明末态光子能量与质心系散射角的关系为
\begin{equation}
  y \approx \frac{1}{1 + (m^2/s)\cdot\text{const}}\ \cdots
    \approx 1 - \frac{m^2}{s}\frac{1}{1-y} .
\end{equation}
于是式~\eqref{eq:5.103} 可以改写为端点附近后向散射光子的能量分布公式：
\begin{equation}
  \frac{d\sigma}{dy} = \frac{2\pi\alpha^2}{s}\ \Bigl[
    \frac{1}{(1-y)+m^2/s}\ \cdots\ + \frac{m^2}{s}\frac{1}{(1-y)^2}
      \frac{1}{(1-y)+m^2/s} \Bigr] ,
  \label{eq:5.104}
\end{equation}
其中方括号中的第一项对应于螺旋度守恒过程，第二项对应于螺旋度翻转过程。因此，例如，如果用一束右手极化的激光束从未极化高能电子束散射，大部分后向散射光子将是右手的，但最高能量的光子将是左手的。这个效应可以在实验上用来测量电子束的极化，或制造具有可调能量分布和极化的高能光子源。

\subsection{对湮灭为光子}

我们仍然可以从康普顿散射振幅中得到一个结果。考虑湮灭过程
\begin{equation}
  e^{+}e^{-} \to 2\gamma ,
\end{equation}
其最低阶由下列图给出
\begin{equation}
  e^{+}(p_2),\ \gamma(k_1),\ \gamma(k_2) \qquad
  \text{with } p \leftrightarrow \text{Compton substitutions} .
\end{equation}
这个过程通过交叉对称性与康普顿散射相联系；我们可以通过作替换从康普顿振幅得到正确的振幅
\begin{equation}
  p \to p_1,\qquad p' \to -p_2,\qquad k \to k_1,\qquad k' \to k_2 .
\end{equation}
在 \eqref{eq:5.87} 中作这些替换，我们得到
\begin{equation}
  \frac{1}{4}\sum_{\text{spins}}|\mathcal{M}|^2 = -2e^4\ \Bigl[
    \frac{p_1\cdot k_1}{p_2\cdot k_1} + \frac{p_1\cdot k_2}{p_2\cdot k_2}
      + 2m^2\Bigl( \frac{1}{p_1\cdot k_1} + \frac{1}{p_1\cdot k_2} \Bigr)
      - m^4\Bigl( \frac{1}{p_1\cdot k_1} + \frac{1}{p_1\cdot k_2} \Bigr)^2
    \Bigr] .
  \label{eq:5.105}
\end{equation}
整体的负号是交叉关系 \eqref{eq:5.68} 的结果，应当去掉。

现在特化到质心系。运动学是标准的，常规计算给出微分截面，
\begin{equation}
  \frac{d\sigma}{d\cos\theta} = \frac{2\pi\alpha^2}{s}\ \Bigl(
    \frac{E}{p} \Bigr)\ \Bigl[
      \frac{E^2 + p^2\cos^2\theta}{m^2 + p^2\sin^2\theta}
      + \frac{2m^2}{m^2+p^2\sin^2\theta}
      - \frac{2m^4}{(m^2+p^2\sin^2\theta)^2}
    \Bigr] .
  \label{eq:5.106}
\end{equation}
在高能极限下，这变为
\begin{equation}
  \frac{d\sigma}{d\cos\theta} \xrightarrow{E\gg m}
    \frac{2\pi\alpha^2}{s}\ \Bigl( \frac{1+\cos^2\theta}{\sin^2\theta}
      \Bigr) .
  \label{eq:5.107}
\end{equation}
HRS 合作组测量的 $E_{\mathrm{cm}} = 29$ GeV 处 $e^{+}e^{-}\to 2\gamma$ 截面的角依赖关系如图~\ref{fig:5.7} 所示；实线是最低阶理论预言，即式~\eqref{eq:5.107}。

\begin{figure}[t]
  \centering
  \includegraphics[width=0.6\textwidth]{fig5.7.png}
  \caption{HRS 合作组（M.~Derrick 等人，\emph{Phys.\ Rev.\ D} \textbf{34}，3286 (1986)）测量的 $E_{\mathrm{cm}} = 29$ GeV 处 $e^{+}e^{-}\to 2\gamma$ 截面的角依赖关系。实线是最低阶理论预言，即式~\eqref{eq:5.107}。}
  \label{fig:5.7}
\end{figure}

\problems

\textbf{5.1} 计算电子-缪子散射 $e^{-}\mu^{-}\to e^{-}\mu^{-}$ 在缪子静止系中的截面，保留电子质量但令 $m_\mu\to\infty$。

\textbf{5.2} \emph{巴巴散射。}计算巴巴散射 $e^{+}e^{-}\to e^{+}e^{-}$ 的微分截面 $d\sigma/d\cos\theta$。你可以取极限 $E_{\mathrm{cm}} \gg m_e$，在此极限下可以忽略电子质量。有两个费曼图；在平方之前必须在不变矩阵元中把它们相加。务必确保这两个图之间的相对符号正确。中间步骤很复杂，但最终结果相当简单。特别是，你可能会发现引入曼德尔施塔姆变量 $s$、$t$ 和 $u$ 是有用的。注意，如果我们忽略电子质量，$s+t+u = 0$。你应该能把微分截面写成如下形式
\begin{equation}
  \frac{d\sigma}{d\cos\theta} = \frac{\pi\alpha^2}{2s}\ \Bigl[ \frac{s^2+u^2}
    {t^2} + \frac{s^2+t^2}{u^2} + \frac{4s^2}{tu} \Bigr] .
\end{equation}
把这个公式改写成用 $\cos\theta$ 表示并画出图形。图中的什么特征导致微分截面在 $\theta\to 0$ 时发散？

\textbf{5.3} 习题3.3中引入的旋量乘积形式体系为计算涉及无质量粒子的树图提供了有效的方法。回忆在习题3.3中我们如下定义旋量乘积：设 $u_{L0}$、$u_{R0}$ 是某个固定的类光动量 $k_0$ 处的左手和右手旋量。它们满足
\begin{equation}
  \slashed{k}_0 u_{L0} = \slashed{k}_0 u_{R0} = 0 .
  \label{eq:5.p1}
\end{equation}
（这些关系只是更标准公式 $\sum u_0\bar{u}_0 = \slashed{k}_0$ 到确定螺旋度上的投影。）然后对任何其他类光动量 $p$ 定义旋量
\begin{equation}
  u_L(p) = \frac{1}{\sqrt{2p\cdot k_0}}\,\slashed{p}\,u_{R0};\qquad
  u_R(p) = \frac{1}{\sqrt{2p\cdot k_0}}\,\slashed{p}\,u_{L0} .
  \label{eq:5.p2}
\end{equation}
我们证明了这些旋量满足 $\slashed{p}u_L(p) = \slashed{p}u_R(p) = 0$；由于周围没有 $m$，它们可以用作费米子或反费米子的旋量。我们定义了
\begin{equation}
  s(p_1, p_2) = \bar{u}_R(p_1)u_L(p_2),\qquad
  t(p_1, p_2) = \bar{u}_L(p_1)u_R(p_2),
\end{equation}
并且，在一个特殊的参考系中，我们证明了性质
\begin{equation}
  t(p_1,p_2) = (s(p_2,p_1))^{*},\qquad
  s(p_1,p_2) = -s(p_2,p_1),\qquad
  |s(p_1,p_2)|^2 = 2p_1\cdot p_2 .
  \label{eq:5.p3}
\end{equation}
现在让我们应用这些结果。
\begin{enumerate}[label=(\alph*)]
  \item 作为热身，给出式~\eqref{eq:5.p3} 中最后一个关系的另一个证明：用 \eqref{eq:5.p1} 把 $|s(p_1,p_2)|^2$ 改写为 Dirac 矩阵的迹，然后应用求迹计算。
  \item 证明，对于任意一串 Dirac 矩阵，
    \begin{equation}
      \slashed{a}\slashed{b}\cdots = \cdots
    \end{equation}
    其中 $\mu, \nu, \rho, \ldots = 0, 1, 2, 3$ 或 5。用这个恒等式证明
    \begin{equation}
      \bar{u}_L(p_1)\gamma^\mu u_L(p_2)
        = \bar{u}_R(p_2)\gamma^\mu u_R(p_1) .
    \end{equation}
  \item 证明费尔兹恒等式
    \begin{equation}
      \bar{u}_L(p_1)\gamma^\mu u_L(p_2)\,\gamma_\mu = 2\bigl[
        u_L(p_2)\bar{u}_L(p_1) + u_R(p_1)\bar{u}_R(p_2) \bigr],
    \end{equation}
    其中 $a, b = 1, 2, 3, 4$ 是 Dirac 指标。这可以通过论证以下陈述来完成：这个方程的右边是一个 Dirac 矩阵；因此，它可以写成第3.4节讨论的 16 个 $\Gamma$ 矩阵的线性组合。它满足
    \begin{equation}
      \gamma^5[M] = -[M]\gamma^5,
    \end{equation}
    因此，它必定具有形式 $V_\mu\gamma^\mu + W_\mu\gamma^\mu\gamma^5$，其中 $V^\mu$ 和 $W^\mu$ 是四矢量。这些四矢量可以用求迹技巧计算；例如，
    \begin{equation}
      V^\mu = \tfrac{1}{4}\mathrm{tr}[M\gamma^\mu] .
    \end{equation}
  \item 考虑过程 $e^{+}e^{-}\to\mu^{+}\mu^{-}$，在 $\alpha$ 的最低阶，忽略电子和缪子两者的质量。首先考虑电子和末态缪子都是右手、正电子和末态反缪子都是左手的情形。（对反缪子使用旋量 $v_R$，对正电子使用 $\bar{u}_R$。）应用费尔兹恒等式证明振幅可以直接用旋量乘积求值。对振幅平方并重现下述结果
    \begin{equation}
      \frac{d\sigma}{d\cos\theta}(e^{-}_R e^{+}_L \to \mu^{-}_R\mu^{+}_L)
    \end{equation}
    式~\eqref{eq:5.21} 中给出的结果。计算该过程的其他螺旋度截面，并证明它们也重现了第5.2节的结果。
  \item 用旋量乘积形式体系，逐螺旋度态地计算无质量电子的巴巴散射微分截面。对初态螺旋度求平均、对末态螺旋度求和，应当重现习题5.2的结果。在这个过程中，你应该能看到这个结果如何作为确定螺旋度贡献之和而出现。
\end{enumerate}

\textbf{5.4} \emph{电子偶素（正电子素）寿命。}
\begin{enumerate}[label=(\alph*)]
  \item 计算极端非相对论极限下 $e^{+}e^{-}$ 湮灭为 2 个光子的振幅 $\mathcal{M}$（即只保留正比于电子和正电子三维动量零次幂的项）。利用这一结果，连同我们关于费米子-反费米子束缚态的形式体系，计算电子偶素 $1S$ 态湮灭为 2 个光子的速率。你应该会发现电子偶素的自旋为 1 的态不会湮灭为 2 个光子，从而证实习题3.8的对称性论证。对于电子偶素自旋为 0 的态，你应该会得到一个正比于 $1S$ 波函数在原点处平方的结果。代入来自非相对论量子力学的这个波函数的值，你应该会得到
    \begin{equation}
      \Gamma = 8.03\times10^9\ \mathrm{sec}^{-1} .
    \end{equation}
    最近的一次测量%
    \footnote{D.~W.~Gidley et.\ al., \emph{Phys.\ Rev.\ Lett.} \textbf{49}, 525
    (1982).} 给出 $\Gamma = 7.994 \pm 0.011\ \mathrm{nsec}^{-1}$；0.5\% 的偏差由辐射修正解释。
  \item 计算更高 $\ell$ 的电子偶素态的衰变率要困难一些；在本习题的其余部分，我们将考虑 $\ell = 1$ 的情形。首先，求出 $e^{+}e^{-}\to 2\gamma$ 振幅中正比于三维动量一次幂的项。（为简单起见，在质心系中计算。）由于
    \begin{equation}
      \psi(\mathbf{x}=0) \sim \int d^3p\ \tilde{\psi}(\mathbf{p}),
    \end{equation}
    振幅的这一部分与 $P$ 波束缚态有重叠。证明 $S=1$ 而不是 $S=0$ 的态可以衰变到 2 个光子。同样，这是 $C$ 的一个推论。
  \item 为了计算这些 $P$ 波态的衰变率，我们需要适当归一化的态矢量。用下式表示三个 $P$ 态波函数
    \begin{equation}
      \psi_i = \xi_i f(|\mathbf{x}|),\qquad \xi_i = \delta^{ij},
    \end{equation}
    归一化为 $\int d^3x\ |\psi_i|^2 = 1$，其 Fourier 变换记为 $\tilde{\psi}_i(\mathbf{p})$。证明
    \begin{equation}
      |B(\mathbf{k})\rangle = \sqrt{2M}\int \frac{d^3p}{(2\pi)^3}\
        \tilde{\psi}_i(\mathbf{p})\ \Xi^i\ a^{\dagger}_{\mathbf{k}/2+\mathbf{p}}
        b^{\dagger}_{-\mathbf{k}/2-\mathbf{p}}|0\rangle
    \end{equation}
    如果 $\Xi^i$ 表示一组归一化为 $\mathrm{tr}[\Xi^{i\dagger}\Xi^j] = \delta^{ij}$ 的三个 $2\times2$ 矩阵，则这是一个适当归一化的束缚态矢量。为了构造 $S=1$ 的态，我们应该让每个 $\Xi^i$ 包含一个 Pauli sigma 矩阵。一般来说，自旋-轨道耦合将按照总角动量 $J$ 分裂 $S=1$、$L=1$ 态的多重态。确定 $J$ 的态由下式给出
    \begin{equation}
      J = 0:\quad \Xi^i = \frac{1}{\sqrt{3}}\sigma^i;\qquad
      J = 1:\quad \Xi^i = -\frac{i}{\sqrt{2}}\epsilon^{ijk}\sigma^k;\qquad
      J = 2:\quad \Xi^i = \frac{1}{\sqrt{2}}h^{ij}\sigma^j,
    \end{equation}
    其中 $\mathbf{n}$ 是满足 $|\mathbf{n}|^2 = 1$ 的极化矢量，$h$ 是无迹张量，其典型值可以是 $h^{12} = 1$ 而所有其他分量为零。
  \item 利用(b)部分导出的 $e^{+}e^{-}\to 2\gamma$ 振幅的展开形式以及(c)部分得到的 $S=1$、$L=1$、确定 $J$ 的电子偶素态的显式形式，对每个 $J$ 计算该态衰变为两个光子的衰变率。
\end{enumerate}
